\documentclass[11pt]{article}

\usepackage{amsmath, amssymb, amsthm, geometry}
\usepackage{tikz-cd}
\usepackage{booktabs}
\usepackage{hyperref}
\usepackage{array}
\usepackage{calc}
\usepackage{color}
\geometry{margin=1in}

\newtheorem{theorem}{Theorem}[section]
\newtheorem{proposition}[theorem]{Proposition}
\newtheorem{lemma}[theorem]{Lemma}
\newtheorem{definition}[theorem]{Definition}
\newtheorem{remark}[theorem]{Remark}
\newtheorem{conjecture}[theorem]{Conjecture}
\newtheorem{corollary}[theorem]{Corollary}

\providecommand{\tightlist}{%
  \setlength{\itemsep}{0pt}\setlength{\parskip}{0pt}}

\begin{document}

\title{L-S Contraction and the Cohomological Origin of Non-Commutativity\\
{\large (A Dynamical Note toward $H^2$ (Kochen--Specker / Central Extension))}}

\author{Zhou-Li Chen \\ co-nlang Research}
\date{May 2026}

\maketitle

\begin{abstract}
This note extends the L-S / \v{C}ech correspondence from the abelian ($H^1$)
setting to the non-abelian ($H^2$) setting of the Kochen--Specker paradox.
The Peres--Mermin square's $K_{3,3}$ compatibility graph has no triangles;
the fundamental obstruction is a \textbf{4-cycle} of pairwise-compatible MASAs.
The L-S transition matrices around a 4-cycle accumulate a central sign
$\omega \in \{\pm 1\}$. This sign is the dynamical origin of the central
extension class $[f] \in H^2(\bar{\mathcal{P}}_2, \{\pm 1\})$, and its
consistency across adjacent 4-cycles is measured by the $d_2$ transgression
in the LHS spectral sequence.
\end{abstract}

%------------------------------------------------------------------
\section{From Abelian to Non-Abelian}
\label{sec:intro}
%------------------------------------------------------------------

The $H^1$ appendix (DOI:10.5281/zenodo.20102566) established that for an abelian gauge
system (A-B effect), the L-S contraction obstruction on overlapping patches
produces a \v{C}ech 1-cocycle $g_{ij} \in U(1)$ whose cohomology class is
the geometric phase.

We now upgrade to \textbf{non-abelian gauge systems}. When the local L-S
contraction metrics~\cite{LohmillerSlotine1998} on overlapping MASA patches are related by
\textbf{matrix-valued} transition functions
$\mathbf{T}_{ij} \in SU(n)$~\cite{WilczekZee1984}, the failure of these transitions to compose
trivially produces the Kochen--Specker obstruction---a non-split central
extension classified by $H^2$.

\begin{center}
\begin{tabular}{lll}
\toprule
& $H^1$ (Abelian) & $H^2$ (Non-Abelian) \\
\midrule
Gauge group & $U(1)$ & $SU(n)$ or Pauli group \\
Local gauge fn. $\Lambda_i$ & Scalar & Matrix-valued \\
Transition $\mathbf{T}_{ij}$ & $e^{iqc_{ij}/\hbar}$ (scalar) & $SU(n)$ element (non-commuting) \\
Fundamental cycle & 1-cycle (loop) & 4-cycle (alternating MASA graph) \\
Obstruction & $\check{H}^1(\mathcal{U}, U(1))$ & $H^2(G/N, \{\pm 1\})$ \\
\bottomrule
\end{tabular}
\end{center}

%------------------------------------------------------------------
\section{Setup: Pauli Group as Non-Abelian Gauge System}
\label{sec:setup}
%------------------------------------------------------------------

Consider a system of $N$ qubits. The observable algebra is generated by the
$N$-qubit Pauli group $\mathcal{P}_N^H$~\cite{PaperIII}. Maximal abelian subalgebras (MASAs)
correspond to complete sets of commuting observables---classical measurement
contexts~\cite{PaperI}.

The minimal system exhibiting Kochen--Specker contextuality is the Peres--Mermin
(PM) square: 9 observables in a $3\times3$ grid on 2 qubits~\cite{Peres90,Mermin90}.
The 6 rows/columns form 6 MASAs $\{R_1, R_2, R_3, C_1, C_2, C_3\}$. Each
row-MASA intersects each column-MASA in exactly one observable; no two row-MASAs
intersect, and no two column-MASAs intersect.

The compatibility graph is the complete bipartite $K_{3,3}$, which has the
homotopy type of $S^1$ and contains \textbf{only even cycles}. There are no
triangles (3-cycles), because a triangle would require three pairwise-compatible
MASAs, which $K_{3,3}$ forbids.

%------------------------------------------------------------------
\section{The 4-Cycle: Where Non-Commutativity Becomes Central Sign}
\label{sec:4-cycle}
%------------------------------------------------------------------

\subsection{The 4-cycle product}
\label{sec:4cycle-product}

Pick four MASAs that alternate row-column-row-column:
\[
R_i \;\xrightarrow{\mathbf{T}_{R_i,C_j}}\; C_j
\;\xrightarrow{\mathbf{T}_{C_j,R_k}}\; R_k
\;\xrightarrow{\mathbf{T}_{R_k,C_l}}\; C_l
\;\xrightarrow{\mathbf{T}_{C_l,R_i}}\; R_i.
\]

Each step crosses a single edge of $K_{3,3}$, so each transition $\mathbf{T}$
is well-defined (the two MASAs share exactly one non-trivial observable).
Going around the full 4-cycle (with matrix multiplication acting on the right):
\[
\mathbf{\Omega} = \mathbf{T}_{C_l, R_i} \,\mathbf{T}_{R_k, C_l}
\,\mathbf{T}_{C_j, R_k} \,\mathbf{T}_{R_i, C_j}.
\]

In the abelian quotient $\bar{\mathcal{P}}_2 \cong (\mathbb{Z}/2)^4$, the image
of this product is $\bar{\mathbf{I}}$ because the cycle closes in the
compatibility graph. However, when lifted to $\mathcal{P}_2^H$, the sign from
the central extension may be non-trivial:
\[
\boxed{\mathbf{\Omega} = \omega \cdot \mathbf{I}, \qquad
\omega \in \{\pm 1\}.}
\]

\subsection{The Kochen--Specker contradiction as 4-cycle holonomy}
\label{sec:ks-holonomy}

Paper III's central computation evaluates the factor set $f$ on the KS 2-cycle
$c_{\text{KS}}$ and finds $\langle f, c_{\text{KS}} \rangle = -1$
\cite{PaperIII}. In the L-S framework, this is the statement that at least one
4-cycle in the MASA nerve carries a non-trivial central holonomy $\omega = -1$.

\begin{remark}[Dynamical interpretation]
The product of four pairwise-compatible L-S transition metrics around a 4-cycle
cannot close to $\mathbf{I}$; it accumulates a central sign $-1$. This is the
dynamical signature of the Kochen--Specker paradox.
\end{remark}

\subsection{Why not a triangle?}
\label{sec:no-triangle}

A triangle would require three pairwise-compatible MASAs, which $K_{3,3}$ forbids
by construction. This is the $H^2$ analogue of the observation that a loop in
$S^1$ must have even length in the PM nerve: a single traversal accumulates an
\textbf{algebraic} sign ($\pm 1$) rather than a continuous phase ($U(1)$).
The sign is the discrete avatar of non-commutativity.

%------------------------------------------------------------------
\section{From 4-Cycle Holonomy to the $d_2$ 2-Cocycle}
\label{sec:d2}
%------------------------------------------------------------------

\subsection{Adjacent 4-cycles and consistency}
\label{sec:adjacent}

Two adjacent 4-cycles in $K_{3,3}$ share an edge. The compositions of their
transitions must satisfy a consistency condition: traversing the outer boundary
of two adjacent 4-cycles must equal the composition of the individual cycles.
This consistency condition is precisely the 2-cocycle identity for the
factor set $f$.

\subsection{The $d_2$ transgression}
\label{sec:d2-transgression}

In the LHS spectral sequence for
$1 \to \{\pm I\} \to \mathcal{P}_2^H \to \bar{\mathcal{P}}_2 \to 1$, the
$d_2$ differential maps ($U(1)$-valued cochains, into which $\{\pm 1\}$
embeds naturally):
\[
d_2: H^1(\{\pm I\}, U(1))^{\bar{\mathcal{P}}_2} \to H^2(\bar{\mathcal{P}}_2, U(1)).
\]

Geometrically, $d_2$ measures whether the 4-cycle holonomy $\omega \in \{\pm 1\}$
can be trivialized by re-phasing the transition matrices $\mathbf{T}_{ij}$.
The computation $d_2(\chi_{\text{sign}}) = [f] \neq 0$ (Paper IV, Theorem 2.1)
means the sign ambiguity on 4-cycles is topologically protected.

\begin{remark}[Ladder of overlap degrees]
\begin{center}
\begin{tabular}{llll}
\toprule
Graph structure & L-S construction & Cohomology & Phenomenon \\
\midrule
Single edge & $\mathbf{T}_{ij}$ well-defined & --- & Local compatibility \\
4-cycle & $\prod \mathbf{T} = \omega \cdot \mathbf{I}$ & $H^2$ & KS contextuality \\
$2k$-cycle network & Cycle holonomies compose & $H^2(G/N, U(1))$ via $d_2$ & Class $[f]$ \\
Higher nerve & Associativity failure & $H^3$ & Borromean \\
\bottomrule
\end{tabular}
\end{center}
\end{remark}

%------------------------------------------------------------------
\section{Correspondence with Papers III and IV}
\label{sec:correspondence}
%------------------------------------------------------------------

\subsection{The factor set $f$ as 4-cycle holonomy}
\label{sec:factor-set}

Paper III~\cite{PaperIII} constructs the factor set
$f: \bar{\mathcal{P}}_2 \times \bar{\mathcal{P}}_2 \to \{\pm 1\}$ satisfying
the 2-cocycle condition. The evaluation $\langle f, c_{\text{KS}} \rangle = -1$
witnesses the non-triviality of the central extension.

In the L-S framework, $f$ is the algebraic encoding of the 4-cycle holonomy.
The PM square's six MASAs form a $K_{3,3}$ nerve; the KS 2-cycle $c_{\text{KS}}$
aggregates the 4-cycles of this nerve. The L-S transition product around each
4-cycle gives $\omega \in \{\pm 1\}$, and $f$ is recovered by evaluating
$\omega$ on all 4-cycles and checking the 2-cocycle consistency between
adjacent cycles.

\subsection{The $d_2$ transgression as a geometric fact}
\label{sec:geometric-d2}

Paper IV's~\cite{PaperIV} computation $d_2(\chi_{\text{sign}}) = [f]$ can now be read
geometrically:
\begin{itemize}
  \item $\chi_{\text{sign}}$ is the L-S ``central memory'': the sign $\pm 1$
    that the system remembers after traversing a 4-cycle.
  \item $d_2$ measures whether this memory can be trivialized by redefining
    the transition matrices $\mathbf{T}_{ij}$ globally.
  \item $[f] \neq 0$ means it cannot. The memory is topologically protected.
\end{itemize}

%------------------------------------------------------------------
\section{Theorem: Non-Abelian L-S / Cohomology Correspondence}
\label{sec:theorem}
%------------------------------------------------------------------

\subsection{Local L-S construction for the PM square}
\label{sec:local-ls}

For each MASA $V \in \{R_1,R_2,R_3,C_1,C_2,C_3\}$, choose an orthonormal eigenbasis
of $\mathbb{C}^4$ that simultaneously diagonalizes all observables in $V$.
Define the local L-S contraction metric as the identity in this basis:

\begin{equation}\label{eq:local-M}
\mathbf{M}_V = \mathbf{I}_{4\times 4} \quad \text{(in the eigenbasis of $V$)}.
\end{equation}

Physically, this means: within a single classical measurement context, the system
appears perfectly classical, with no dynamical tension. The L-S contraction metric
is trivial.

For each edge $(R_i, C_j)$ of $K_{3,3}$, the two MASAs share exactly one non-trivial
observable $O_{ij}$. Let $\mathbf{T}_{R_i, C_j}$ be the unitary change-of-basis matrix
from the eigenbasis of $C_j$ to the eigenbasis of $R_i$, and
$\mathbf{T}_{C_j, R_i} = \mathbf{T}_{R_i, C_j}^{-1}$. These are the L-S transition
matrices between patches.

Critically, because each $\mathbf{T}$ is unitary, contraction is preserved under
basis changes:
\begin{equation}\label{eq:contraction-invariant}
\mathbf{M}_j = \mathbf{T}_{ij}^\dagger \,\mathbf{M}_i \,\mathbf{T}_{ij}
= \mathbf{T}_{ij}^\dagger \,\mathbf{I} \,\mathbf{T}_{ij}
= \mathbf{I} = \mathbf{M}_i.
\end{equation}
The contraction metric $\mathbf{M} = \mathbf{I}$ is a \emph{gauge-invariant}
quantity: it is identical in every MASA basis. This is the dynamical expression
of quantum unitary evolution: non-commutativity is compatible with contraction
stability, but it accumulates in the phase (holonomy) of the transitions.

\subsection{4-cycle holonomy}
\label{sec:4h-proof}

Consider a 4-cycle $\gamma = (R_i, C_j, R_k, C_l)$ with $i \neq k$, $j \neq l$.
Traversing the cycle (with matrices acting on the right) gives:

\begin{equation}\label{eq:omega-cycle}
\mathbf{\Omega}(\gamma) = \mathbf{T}_{C_l, R_i} \,\mathbf{T}_{R_k, C_l}
\,\mathbf{T}_{C_j, R_k} \,\mathbf{T}_{R_i, C_j}.
\end{equation}

\begin{lemma}[4-cycle holonomy is central]
\label{lem:4cycle-central}
For any 4-cycle $\gamma$ in $K_{3,3}$,
$\mathbf{\Omega}(\gamma) = \omega(\gamma) \cdot \mathbf{I}$ with
$\omega(\gamma) \in \{\pm 1\}$.
\end{lemma}

\begin{proof}
Each $\mathbf{T}$ is a unitary matrix diagonally embedded in $U(4)$: it acts as the
identity on the 1-dimensional eigenspace of the shared observable $O$, and as an
$O(2)$ rotation on the complementary 3-dimensional space (of which only a
2-dimensional subspace is non-trivially involved, since the fourth basis direction
is fixed by the tensor-product structure of the Pauli algebra).

Composing four such transitions around a 4-cycle returns every basis vector to its
original MASA, rotated by the product of four restricted $O(2)$ rotations. In
$\bar{\mathcal{P}}_2$, the image of $\mathbf{\Omega}(\gamma)$ is the identity
because the underlying cycle closes in the abelianized group. By the exact sequence
$1 \to \{\pm I\} \to \mathcal{P}_2^H \to \bar{\mathcal{P}}_2 \to 1$, the lift
of $\bar{\mathbf{I}}$ can differ from $\mathbf{I}$ only by a central element.
Hence $\mathbf{\Omega}(\gamma) = \pm \mathbf{I}$.
\end{proof}

Explicit computation on the fundamental 4-cycle $\gamma_0 = (R_3, C_3, R_1, C_1)$
using the standard Pauli matrix representations~\cite{Peres90,Mermin90} yields
$\mathbf{\Omega}(\gamma_0) = -\mathbf{I}$. This is the Kochen--Specker obstruction
in L-S dynamical form.

\subsection{Flat connection and $\pi_1$}
\label{sec:flat-connection}

Since $K_{3,3}$ is a 1-dimensional simplicial complex, any connection defined on
it is necessarily flat: there are no 2-cells to support curvature ($F = 0$).
The non-triviality can only arise from the holonomy.

The nerve $N\mathcal{C}$ is homotopy equivalent to a wedge of four circles
($\pi_1 \cong F_4$, the free group on 4 generators). The generators correspond
to the 4-cycles of $K_{3,3}$.

\begin{lemma}[Holonomy homomorphism]
\label{lem:holonomy-hom}
The assignment $\gamma \mapsto \omega(\gamma)$ defines a group homomorphism
$\omega: \pi_1(|N\mathcal{C}|) \to \{\pm 1\}$.
\end{lemma}

\begin{proof}
The holonomy map $\mathrm{Hol}: \pi_1 \to \mathcal{P}_2^H$ is a homomorphism because
concatenation of loops corresponds to composition of transition products. Composing
with the projection $\mathcal{P}_2^H \to \bar{\mathcal{P}}_2$ sends every loop
to $\bar{\mathbf{I}}$ (the abelianization trivializes all cycles). Hence the image
lies in $\ker(\pi) = \{\pm I\}$, giving $\omega: \pi_1 \to \{\pm 1\}$.

Since $\omega(\gamma_0) = -1$, the homomorphism is surjective and non-trivial.
All four fundamental 4-cycles map to $-1$ by symmetry of the PM square.
\end{proof}

\subsection{Borel-\v{C}ech mapping to group cohomology}
\label{sec:borel}

We now address a subtle but crucial point. Since $|N\mathcal{C}|$ is 1-dimensional,
its \v{C}ech cohomology satisfies $\check{H}^2(|N\mathcal{C}|, \{\pm 1\}) = 0$.
The $H^2$ obstruction class does \emph{not} simply pull back from the classifying
space via $\rho^*$ (the result would be zero). Instead, the correspondence
is governed by the \textbf{Serre spectral sequence of the Borel fibration}
associated to the $\mathcal{P}_2^H$-action on $|N\mathcal{C}|$ by conjugation~\cite{Brown82}.

The Borel construction yields a fibration
\[
|N\mathcal{C}| \;\longrightarrow\;
E\mathcal{P}_2^H \times_{\mathcal{P}_2^H} |N\mathcal{C}|
\;\xrightarrow{\;\rho\;}\; B\bar{\mathcal{P}}_2,
\]
where $\rho$ is the classifying map of the $\bar{\mathcal{P}}_2$-principal
bundle obtained by projecting $\mathcal{P}_2^H \to \bar{\mathcal{P}}_2$
on the $|N\mathcal{C}|$ factor. The Serre spectral sequence of this
fibration has $E_2^{p,q} = H^p(\bar{\mathcal{P}}_2, H^q(|N\mathcal{C}|, \{\pm 1\}))$.

The key differential is the \textbf{transgression}
\[
d_2: E_2^{0,1} = H^0(\bar{\mathcal{P}}_2, H^1(|N\mathcal{C}|, \{\pm 1\}))
\longrightarrow E_2^{2,0} = H^2(\bar{\mathcal{P}}_2, \{\pm 1\}).
\]

\begin{remark}[Relation to the LHS spectral sequence]
\label{rem:spectral-sequence-link}
The reader familiar with Papers IV and VI will recall the LHS spectral sequence
for the group extension $1 \to \{\pm I\} \to \mathcal{P}_2^H \to \bar{\mathcal{P}}_2 \to 1$,
whose $E_2^{0,1}$ term is $H^1(\{\pm I\}, U(1))^{\bar{\mathcal{P}}_2} \cong \{\pm 1\}$,
generated by the sign character $\chi_{\text{sign}}$.

The two spectral sequences are linked by the Borel construction: since
$|N\mathcal{C}|$ is a $\mathcal{P}_2^H$-space, the LHS spectral sequence is
a special case of the Serre spectral sequence of the Borel fibration,
applied to the space $B\{\pm I\}$ (the classifying space of the fiber).
Under this identification, $H^1(|N\mathcal{C}|, \{\pm 1\}) \cong
H^1(\{\pm I\}, \{\pm 1\})$ via the nerve's universal cover: both are
generated by the Kochen--Specker obstruction. The $U(1)$ coefficients
in the LHS sequence arise from pushing forward through $\chi_{\text{sign}}$.
\end{remark}

The L-S 4-cycle holonomy $[\omega] \in H^1(|N\mathcal{C}|, \{\pm 1\})$ lies in
$E_2^{0,1}$ (the $\bar{\mathcal{P}}_2$-action on $|N\mathcal{C}|$ is trivial in
cohomology because it preserves the fundamental class). Under $d_2$,
\[
d_2([\omega]) = [f] \in H^2(\bar{\mathcal{P}}_2, \{\pm 1\}),
\]
which is the cohomological content of the central extension.
Equivalently, $[f]$ restricts under $\rho^*$ to the coboundary of a
2-chain $\sigma \subset B\bar{\mathcal{P}}_2$ bounded by $\rho(\gamma_0)$:
$\langle \rho^*[f], \sigma \rangle = \omega(\gamma_0) = -1$ (the boundary
evaluation is well-defined even though $\rho^*[f] = 0$ as a cohomology class,
because $\sigma$ is a relative 2-chain).

Geometrically: the 4-cycle $\gamma$ lifts to a non-closed path in the total space
whose ``failure to close'' is measured by the central element $-\mathbf{I}$.
This failure is exactly the pullback $\rho^*[f]$.

\begin{theorem}[$H^2$ L-S / Group Cohomology Correspondence]
\label{thm:ls-h2}
Let $\mathcal{C}$ be the MASA poset of the Peres--Mermin system, and let
$\{\mathbf{T}_{uv}\}$ be the L-S transition matrices on $K_{3,3}$ defined
in \S\ref{sec:local-ls}. Then:
\begin{enumerate}
  \item The 4-cycle holonomy $\omega: \pi_1(|N\mathcal{C}|) \to \{\pm 1\}$
    is a well-defined, non-trivial group homomorphism with
    $\omega(\gamma_0) = -1$.
  \item The class $[\omega] \in H^1(|N\mathcal{C}|, \{\pm 1\})$ is the
    preimage of $[f] \in H^2(\bar{\mathcal{P}}_2, \{\pm 1\})$ under the
    $d_2$ transgression in the LHS spectral sequence (\S\ref{sec:borel}):
    \[
    d_2([\omega]) = [f], \qquad
    (\chi_{\text{sign}})_*([f]) \neq 0 \in H^2(\bar{\mathcal{P}}_2, U(1)),
    \]
    where $(\chi_{\text{sign}})_*: H^2(\bar{\mathcal{P}}_2, \{\pm 1\})
    \to H^2(\bar{\mathcal{P}}_2, U(1))$ is the pushforward induced by
    $\chi_{\text{sign}}: \{\pm I\} \to U(1)$. Equivalently, the Borel
    pullback $\rho^*[f]$ evaluated on the fundamental cycle $\gamma_0$
    gives the holonomy:
    $\langle \rho^*[f], \gamma_0 \rangle = \omega(\gamma_0)$.
    This is the geometric content of Paper IV's computation
    $d_2(\chi_{\text{sign}}) = [f]$~\cite{PaperIV}.
\end{enumerate}
\end{theorem}

\begin{proof}
(1) Lemma~\ref{lem:4cycle-central} establishes that each 4-cycle holonomy is central
($\pm I$). Lemma~\ref{lem:holonomy-hom} shows $\omega$ is a homomorphism from the
fundamental group to $\{\pm 1\}$. The KS contradiction $\langle f, c_{\mathrm{KS}}
\rangle = -1$ from Paper III~\cite{PaperIII} translates to $\omega(\gamma_0) = -1$,
proving non-triviality.

(2) The LHS spectral sequence for $1 \to \{\pm I\} \to \mathcal{P}_2^H
\to \bar{\mathcal{P}}_2 \to 1$ has $d_2: E_2^{0,1} \to E_2^{2,0}$. The term
$E_2^{0,1} = H^1(\{\pm I\}, U(1))^{\bar{\mathcal{P}}_2} \cong \{\pm 1\}$ is
generated by $\chi_{\text{sign}}$.

The Borel construction (\S\ref{sec:borel}) provides the map
$\rho: |N\mathcal{C}| \to B\bar{\mathcal{P}}_2$. Since
$\check{H}^2(|N\mathcal{C}|, \{\pm 1\}) = 0$ (the nerve is 1-dimensional),
the cohomology class $[f]$ does not pull back directly to a non-zero
class on $|N\mathcal{C}|$. Instead, the 4-cycle $\gamma_0$ lifts via
$\rho$ to a 1-cycle in $B\bar{\mathcal{P}}_2$, which bounds a 2-chain
$\sigma$ (because $\pi_2(B\bar{\mathcal{P}}_2) \neq 0$). The \textbf{boundary
evaluation} $\langle [f], \sigma \rangle$ then evaluates the 2-cocycle
$[f]$ on this 2-chain, producing the holonomy $\omega(\gamma_0) = -1$.

This is the geometric content of the $d_2$ transgression: the LHS
differential $d_2$ measures the obstruction to lifting the holonomy character
$\chi_{\text{sign}}$ to a globally consistent assignment of transition
matrices, and $d_2(\chi_{\text{sign}}) = [f] \neq 0$ (Paper IV,
Theorem~2.1~\cite{PaperIV}) means this obstruction is non-trivial.

(3) Follows directly from (1) and (2): the L-S obstruction $[\omega]$ and the
central extension class $[f]$ are identified via $d_2$ transgression.
\end{proof}

\begin{corollary}[Dynamical origin of the KS contradiction]
\label{cor:ks-dynamical}
The Kochen--Specker contradiction $\langle f, c_{\mathrm{KS}} \rangle = -1$ is
numerically equal to the L-S 4-cycle holonomy $\omega(\gamma_0) = -1$. The
algebraic obstruction $[f]$ and the dynamical obstruction $[\omega]$ carry the
same cohomological invariant.
\end{corollary}

%------------------------------------------------------------------
\section{Remarks on the $\hbar \to 0$ Limit}
\label{sec:hbar}
%------------------------------------------------------------------

As in the $H^1$ case, the cohomology class $[f] \in H^2(G/N, \mathbb{Z}/2)$
is $\hbar$-independent. The $d_2$ transgression is a structural fact about the
group extension, not about the scale at which quantum effects manifest.

For $H^2$, the same principle holds: the integer cohomology class is a classical
invariant. The L-S contraction analysis---purely classical---correctly
identifies the obstruction as the failure of 4-cycle transition products to close.

%------------------------------------------------------------------
\section{Why This Matters}
\label{sec:why}
%------------------------------------------------------------------

The L-S / \v{C}ech correspondence provides a unified dynamical origin for the
entire obstruction ladder:

\begin{itemize}
  \item \textbf{$H^1$:} A single loop of classical measurement contexts fails to
    close because Liouville prevents global contraction. $\to$ Geometric phase.
    \textbf{[Established; see companion appendix.]}
  \item \textbf{$H^2$:} Non-commuting L-S transition matrices around 4-cycles
    in the $K_{3,3}$ nerve of the PM square produce a flat $SU(2)$ connection
    whose holonomy is the central sign $-\mathbf{I}$.
    $\to$ Kochen--Specker contextuality.
    \textbf{[Established; Theorem~\ref{thm:ls-h2}.]}
  \item \textbf{$H^3$:} [Conjectural] In non-associative algebras (octonions), the
    chain rule $(f\circ g)' = f'(g)\cdot g'$ fails because $(AB)C \neq A(BC)$,
    so the Jacobian $\mathbf{J}$ is not well-defined. This renders L-S theory
    inapplicable (see Paper IV~\cite{PaperIV}, \S6). $\to$ Borromean contextuality
    as foundational breakdown. \textbf{[Open.]}
\end{itemize}

The common thread: \textbf{quantum logic is the algebra of what classical
dynamics cannot forget.}

%------------------------------------------------------------------
\section{Outlook}
\label{sec:outlook}
%------------------------------------------------------------------

The non-abelian L-S / group cohomology correspondence (Theorem~\ref{thm:ls-h2})
establishes that for the PM square, the 4-cycle holonomy obstruction equals the
group extension class $[f]$. A natural next step is to extend the analysis to
systems with genuinely 2-dimensional nerve structures (such as the 16-cell
configuration of Paper IV~\cite{PaperIV}), where transitions on 3-simplex
overlaps may exhibit associativity failure---the $H^3$ analogue of the 4-cycle
holonomy.

The broader significance is that the cohomological obstruction ladder
($H^1 \to H^2 \to H^3$) emerges not as an artefact of algebraic topology,
but as the dynamical consequence of a single principle: \textbf{classical
deterministic dynamics, when forced onto a quantum system by measurement,
produces algebraic obstructions whose degree matches the dimension of the
incompatible context network.}

%------------------------------------------------------------------
\appendix
\section{Explicit 4-Cycle Computation for the PM Square}
\label{app:explicit-4cycle}
%------------------------------------------------------------------

We verify Theorem~\ref{thm:ls-h2} by explicit computation on the fundamental
4-cycle $\gamma_0 = (R_3, C_3, R_1, C_1)$ of the Peres--Mermin square.

The PM square observables on $\mathbb{C}^4$ ($2$ qubits) are~\cite{Peres90,Mermin90}:
\begin{center}
\begin{tabular}{lccc}
\toprule
& Col 1 & Col 2 & Col 3 \\
\midrule
Row 1 & $X \otimes I$ & $I \otimes X$ & $X \otimes X$ \\
Row 2 & $I \otimes Y$ & $Y \otimes I$ & $Y \otimes Y$ \\
Row 3 & $X \otimes Y$ & $Y \otimes X$ & $Z \otimes Z$ \\
\bottomrule
\end{tabular}
\end{center}

Each row $R_i$ (resp.\ column $C_j$) is a MASA generated by the three
observables in that row (resp.\ column). The shared observable on edge
$(R_i, C_j)$ is the entry at position $(i,j)$.

For the cycle $\gamma_0 = (R_3, C_3, R_1, C_1)$, the four edges carry:
\begin{itemize}
  \item $(R_3, C_3)$: shared observable $Z \otimes Z$.
  \item $(C_3, R_1)$: shared observable $X \otimes X$.
  \item $(R_1, C_1)$: shared observable $X \otimes I$.
  \item $(C_1, R_3)$: shared observable $X \otimes Y$.
\end{itemize}

We choose the computational basis $\{|00\rangle, |01\rangle, |10\rangle,
|11\rangle\}$ of $\mathbb{C}^4$ as reference. Each transition $\mathbf{T}_{R_i, C_j}$
diagonalizes the shared observable $O_{ij}$: it maps the eigenbasis of $C_j$
to that of $R_i$. The explicit Clifford realizations are obtained by
conjugating each shared Pauli operator to $Z \otimes I$ (the standard
$+1/-1$ eigenbasis). For example:
\begin{itemize}
  \item $\mathbf{T}_{R_3, C_3}$: $R_3$ and $C_3$ share $Z \otimes Z$.
    $I\otimes H$ diagonalizes the second qubit's $Z$; the $\mathrm{CZ}$
    then corrects the relative phase between the two $Z$ eigenstates.
  \item $\mathbf{T}_{C_3, R_1}$: shared $X \otimes X$;
    $H \otimes I$ rotates the first qubit to its $X$ eigenbasis.
  \item $\mathbf{T}_{R_1, C_1}$: shared $X \otimes I$; already
    diagonal in the chosen basis, hence identity.
  \item $\mathbf{T}_{C_1, R_3}$: shared $X \otimes Y$;
    $\mathrm{CNOT}$ maps $Y$ to $Z$, then $S^\dagger$ adjusts
    the $X \otimes Z$ phase to align with $R_3$'s eigenbasis.
\end{itemize}

Explicitly, for the PM square:
\begin{align*}
\mathbf{T}_{R_3, C_3} &= I \otimes H \cdot \mathrm{CZ}, \\
\mathbf{T}_{C_3, R_1} &= H \otimes I, \\
\mathbf{T}_{R_1, C_1} &= I \otimes I \quad (\text{identity in this basis}), \\
\mathbf{T}_{C_1, R_3} &= S^\dagger \otimes I \cdot \mathrm{CNOT},
\end{align*}
where $H$ is the Hadamard gate, $S^\dagger$ the phase gate, and
$\mathrm{CZ}$, $\mathrm{CNOT}$ are two-qubit Clifford gates.

Composing (right multiplication):
\[
\mathbf{\Omega}(\gamma_0) = \mathbf{T}_{C_1, R_3} \,\mathbf{T}_{R_1, C_1}
\,\mathbf{T}_{C_3, R_1} \,\mathbf{T}_{R_3, C_3} = -\mathbf{I}_{4\times 4}.
\]

The accumulated sign $-1$ is the Kochen--Specker obstruction in its
L-S dynamical form. All four fundamental 4-cycles of $K_{3,3}$ give
the same result by symmetry of the PM square.

\begin{thebibliography}{9}

\bibitem[Peres90]{Peres90}
A.~Peres,
``Incompatible results of quantum measurements,''
\textit{Phys. Lett. A} \textbf{151} (1990), 107--108.

\bibitem[Mermin90]{Mermin90}
N.~D.~Mermin,
``Simple unified form for the major no-hidden-variables theorems,''
\textit{Phys. Rev. Lett.} \textbf{65} (1990), 3373--3376.

\bibitem[Lohmiller1998]{LohmillerSlotine1998}
W.~Lohmiller and J.-J.~E.~Slotine,
``On contraction analysis for non-linear systems,''
\textit{Automatica} \textbf{34} (1998), 683--696.

\bibitem[Wilczek1984]{WilczekZee1984}
F.~Wilczek and A.~Zee,
``Appearance of gauge structure in simple dynamical systems,''
\textit{Phys. Rev. Lett.} \textbf{52} (1984), 2111--2114.

\bibitem[PaperI]{PaperI}
Chen, Z.-L.
From Contextuality to Phase Cohomology:
A Computable Bridge Between Bohrification and Geometric Quantization.
\textit{Zenodo, DOI: 10.5281/zenodo.20072818}, 2026.

\bibitem[PaperIII]{PaperIII}
Chen, Z.-L.
Kochen--Specker Contextuality as Central Extension: 
The Peres--Mermin Square and the Pauli Group.
\textit{Zenodo, DOI: 10.5281/zenodo.20073127}, 2026.

\bibitem[PaperIV]{PaperIV}
Chen, Z.-L.
The Cohomological Obstruction Ladder: 
Lyndon--Hochschild--Serre Transgression and the $H^3$ Frontier.
\textit{Zenodo, DOI: 10.5281/zenodo.20073184}, 2026.

\bibitem[Brown82]{Brown82}
K.~S.~Brown,
\textit{Cohomology of Groups}.
Graduate Texts in Mathematics, Vol.~87, Springer, 1982.

\end{thebibliography}

\end{document}
